% intro.tex (v15, 2026-10-21): four paragraphs: why ALT and its observations are the problem, what the two modelling routes and the products leave open, what we do, what the paper delivers.
The active layer is the near-surface ground above permafrost that thaws each summer and refreezes each winter. Its thickness (ALT) sets the depth of seasonal thaw and controls ground stability, near-surface hydrology and the amount of soil organic carbon exposed to decomposition\cite{Hjort2018,Schuur2015,Biskaborn2019}. ALT responds to the thermal conditions of each year, a variation recorded by long-term programmes such as the Circumpolar Active Layer Monitoring (CALM) network\cite{Nelson1997,Streletskiy2025}. The observations, however, are few and clustered: most come from a small number of campaigns and long-term grids, many sites are visited once, and large parts of the permafrost region have none\cite{Moore2025,Veremeeva2025}. Any ALT map therefore rests on a model that extends sparse observations to unobserved cells and years, and the question is which model does so reliably.

Two routes exist. The Stefan solution relates the thaw depth to the thawing degree-days (TDD) of the season through one coefficient that summarises the thermal properties and moisture of the soil\cite{Stefan1891,Riseborough2008,Nelson1997}. It responds to the climate of each year and, having one parameter, can be refitted with few observations and applied far from them; its weakness is that one coefficient per region cannot represent the soils and terrain within the region, which leave systematic residuals. Machine learning (ML) takes the opposite position: it fits the training sites closely from climate, terrain and soil predictors, but it must learn the climate response from the labels, and its accuracy falls at locations unlike the training sites when it is validated with spatially separated data\cite{Roberts2017,Ploton2020,MeyerPebesma2022,Wadoux2021}. The gridded ALT products that now cover the northern permafrost region sit on these two routes. The ESA CCI product solves a ground heat-transfer model driven by satellite land-surface temperature\cite{Westermann2024,Westermann2025}; the ML products of Liu et al.\cite{Liu2024} and Wei et al.\cite{Wei2026} relate site-year observations to yearly predictors, and the products of Aalto et al.\cite{Aalto2018} and Ran et al.\cite{Ran2022a} map a multi-year mean. The ML products are validated by random or leave-one-site-out splits within their own site sets, with reported errors from 21.6 cm\cite{Liu2024} to 60.5 cm\cite{Wei2026} that are not comparable with each other, and none has been evaluated against observations of a stated year in a region that contributed no training sites.

Physics-guided learning, in which a physical model carries the dominant response and a learner corrects what the physics omits, has improved accuracy and extrapolation in lake, river and soil applications\cite{Willard2022,Read2019,Jia2021}, and in Alaska the Stefan model retained its accuracy on held-out sites better than a random forest did\cite{Gautam2025}. Whether a learner that corrects the Stefan model keeps the transferability of the physics while gaining the accuracy of ML has not been tested for ALT. Three points decide its practical value and have not been established: the conditions under which the correction improves on the physics, the number of local observations a new region needs, and how such a model compares with direct ML and with the published products when all are scored on the same sites and years.

Here we address these three points with a site-year design, in which every label, climate input, prediction and score refers to one observation site in one stated year, and the climate input is computed from ERA5-Land for that year and accumulated to the observation day\cite{munozsabater2021_era5_land}. With \nncyAll{} site-year observations of 2010--2024 from Alaska, the Lena Delta and north-western Canada, we compare the Stefan model, direct ML and a physics-residual workflow, in which gradient boosting learns the residual of a locally recalibrated Stefan model from 23 climate, terrain and soil inputs, under five validation designs that range from random splits to the hold-out of an entire region, and we score five published products on the same site-years. The study (i) shows how and why the ranking of the methods changes with the validation design and where direct ML fails, (ii) identifies the conditions and the number of local labels under which the residual learner improves on the physics, (iii) explains the regional differences between the published products, and (iv) delivers year-labelled 1 km ALT maps of the three regions for 2011--2024 with an area of applicability and prediction intervals.
